\documentclass[12pt,a4paper]{article}

% Paquetes
\usepackage[utf8]{inputenc}
\usepackage[spanish]{babel}
\usepackage{amsmath, amssymb}
\usepackage{graphicx}
\usepackage{float}
\usepackage{subcaption}
\usepackage[margin=2.5cm]{geometry}
\usepackage{hyperref}
\usepackage{listings}
\usepackage{xcolor}
\usepackage{circuitikz}
\usepackage{booktabs}
\usepackage{fancyhdr}
\usepackage{titlesec}

% Configuracion de listings para MATLAB
\lstdefinestyle{matlab}{
    language=Matlab,
    basicstyle=\ttfamily\small,
    keywordstyle=\color{blue}\bfseries,
    commentstyle=\color{green!50!black},
    stringstyle=\color{red},
    numbers=left,
    numberstyle=\tiny\color{gray},
    frame=single,
    breaklines=true,
    captionpos=b
}

% Header/Footer
\pagestyle{fancy}
\fancyhf{}
\fancyhead[L]{Identificaci\'on de Sistemas}
\fancyhead[R]{Circuito RC de Primer Orden}
\fancyfoot[C]{\thepage}

% Datos
\title{
    \textbf{Identificaci\'on de Sistemas} \\[0.5cm]
    \Large Circuito RC de Primer Orden: Simulaci\'on y Realizaci\'on F\'isica \\[1cm]
    \large Basado en: \textit{Engineering Circuit Analysis, 8th Ed.} \\
    W. H. Hayt, J. E. Kemmerly, S. M. Durbin \\
    Cap\'itulos 7 y 8
}
\author{Equipo: Malvados y Asociados}
\date{\today}

\begin{document}

\maketitle
\tableofcontents
\newpage

% ============================================================
\section{Introducci\'on}
% ============================================================

% TODO: Describir el objetivo del proyecto y contexto del curso

\subsection{Objetivo}
Simular y analizar la respuesta de un circuito RC de primer orden a un tren de pulsos, identificando los cuatro casos caracter\'isticos de la respuesta transitoria para determinar el valor de la constante de tiempo $\tau$.

\subsection{Alcance}
\begin{enumerate}
    \item \textbf{Fase 1:} Simulaci\'on en MATLAB --- Modelado matem\'atico y simulaci\'on de los 4 casos.
    \item \textbf{Fase 2:} Realizaci\'on f\'isica --- Implementaci\'on con instrumentaci\'on electr\'onica.
\end{enumerate}

% ============================================================
\section{Marco Te\'orico}
% ============================================================

\subsection{Circuito RC de Primer Orden}

% TODO: Completar con base en el libro (Cap. 7-8)

\begin{figure}[H]
    \centering
    \begin{circuitikz}[american]
        \draw (0,0) to[V, v=$v_s(t)$] (0,3)
                    to[R, l=$R$] (4,3)
                    to[C, l=$C$, v=$v_c(t)$] (4,0)
                    -- (0,0);
    \end{circuitikz}
    \caption{Circuito RC serie de primer orden.}
    \label{fig:circuito_rc}
\end{figure}

\subsection{Ecuaci\'on Diferencial}

La ecuaci\'on diferencial que describe el voltaje en el capacitor es:

\begin{equation}
    RC \frac{dv_c(t)}{dt} + v_c(t) = v_s(t)
    \label{eq:edo_rc}
\end{equation}

donde $\tau = RC$ es la constante de tiempo del circuito.

\subsection{Soluci\'on General}

La soluci\'on a la ecuaci\'on diferencial es:

\begin{equation}
    v_c(t) = V_f + (V_i - V_f) e^{-t/\tau}
    \label{eq:solucion_general}
\end{equation}

donde:
\begin{itemize}
    \item $V_i$ = voltaje inicial del capacitor
    \item $V_f$ = voltaje final (estado estable)
    \item $\tau = RC$ = constante de tiempo
\end{itemize}

\subsection{Respuesta a Tren de Pulsos: Los 4 Casos}

% TODO: Documentar los 4 casos basados en el libro

% Caso 1:
% Caso 2:
% Caso 3:
% Caso 4:

% ============================================================
\section{Fase 1: Simulaci\'on en MATLAB}
% ============================================================

\subsection{Metodolog\'ia}

% TODO: Describir el enfoque de simulaci\'on

\subsection{C\'odigo y Resultados}

% TODO: Incluir scripts y gr\'aficas

% \lstinputlisting[style=matlab, caption={Script principal}]{../../scripts/rc_simulation.m}

\subsection{An\'alisis de Resultados}

% TODO: Interpretar los 4 casos

% ============================================================
\section{Fase 2: Realizaci\'on F\'isica}
% ============================================================

% TODO: Segunda fase del proyecto

\subsection{Dise\~no del Circuito}

\subsection{Instrumentaci\'on}

\subsection{Mediciones y Resultados}

\subsection{Comparaci\'on Simulaci\'on vs. Experimental}

% ============================================================
\section{Conclusiones}
% ============================================================

% TODO: Conclusiones finales

% ============================================================
\section{Problemas Resueltos}
% ============================================================

% TODO: Problemas de final de cap\'itulo 7 y 8

% ============================================================
\appendix
\section{C\'odigo MATLAB Completo}
% ============================================================

% TODO: Anexar scripts completos

\end{document}
